\documentclass[11pt]{article}

\usepackage[margin=1in]{geometry}
\usepackage{amsmath,amssymb}
\usepackage{microtype}
\usepackage[hidelinks]{hyperref}
\usepackage{enumitem}

\title{From Two-Digit Multiplication to Convolution\\
\large Positional Numbers as Polynomials in Their Base}
\author{Albert Stockton}
\date{}

\begin{document}

\maketitle

\begin{abstract}
The familiar algorithm for multiplying multidigit numbers can be understood as a form of polynomial multiplication. Beginning with the multiplication of two two-digit decimal numbers, this article develops a general formula for multiplication in any positional number system. A numeral in base $B$ is interpreted as a polynomial evaluated at $B$, while multiplication produces a discrete convolution of the digit sequences. The usual process of carrying is then shown to be a normalization procedure that converts unrestricted polynomial coefficients into valid base-$B$ digits. This perspective connects elementary place-value arithmetic with algebra, number systems, signal processing, and modern algorithms for multiplying large integers.
\end{abstract}

\section{Introduction}

Consider two two-digit decimal numbers whose digits are $a,b,c,$ and $d$:
\[
\overline{ab}=10a+b,
\qquad
\overline{cd}=10c+d.
\]
Their product is
\[
(10a+b)(10c+d).
\]
Applying the distributive property gives
\[
100ac+10ad+10bc+bd,
\]
or equivalently,
\[
\boxed{100ac+10(ad+bc)+bd}.
\]

This formula may initially look like a shortcut for multiplying two-digit numbers. However, it reveals a deeper structure. The powers $100$, $10$, and $1$ are
\[
10^2,\qquad 10^1,\qquad 10^0.
\]
The digits act like polynomial coefficients, while the number $10$ acts like the value assigned to a variable. Once this observation is generalized, ordinary multiplication becomes polynomial multiplication followed by carrying.

The identity itself is established mathematics rather than a new theorem. The aim here is expository: to show how a simple observation about decimal digits naturally leads to positional notation, convolution, and fast multiplication algorithms.

\section{Two-Digit Multiplication}

Let
\[
X=10a+b
\]
and
\[
Y=10c+d,
\]
where $a$ and $c$ are the tens digits and $b$ and $d$ are the ones digits. Then
\begin{align*}
XY
&=(10a+b)(10c+d)\\
&=100ac+10ad+10bc+bd\\
&=100ac+10(ad+bc)+bd.
\end{align*}

For example, take
\[
23\times45.
\]
Here,
\[
a=2,\quad b=3,\quad c=4,\quad d=5.
\]
Therefore,
\begin{align*}
23\times45
&=100(2)(4)+10\bigl((2)(5)+(3)(4)\bigr)+(3)(5)\\
&=800+10(10+12)+15\\
&=800+220+15\\
&=1035.
\end{align*}

The three groups correspond to powers of ten:
\[
ac\,10^2+(ad+bc)10^1+bd\,10^0.
\]
The coefficient of $10^1$ contains two terms because there are two ways to obtain a total exponent of one:
\[
10^1\cdot10^0
\qquad\text{and}\qquad
10^0\cdot10^1.
\]
This is the first appearance of the convolution pattern.

\section{Positional Representation}

In base $B$, a numeral with digits
\[
a_na_{n-1}\cdots a_1a_0
\]
represents the number
\[
\boxed{A=\sum_{i=0}^{n}a_iB^i}.
\]
The digit $a_i$ occupies position $i$, measured from right to left beginning with position zero.

For example, in base ten,
\[
527=5(10^2)+2(10^1)+7(10^0).
\]
In base two,
\[
1011_2=1(2^3)+0(2^2)+1(2^1)+1(2^0)=11_{10}.
\]
In base sixteen,
\[
2A_{16}=2(16^1)+10(16^0)=42_{10}.
\]

For an ordinary nonnegative integer base $B\ge 2$, the allowed digits are
\[
0,1,\ldots,B-1.
\]
Restricting the coefficients to this range gives every nonnegative integer a unique finite base-$B$ representation, apart from unnecessary leading zeros.

The expression
\[
a_nB^n+\cdots+a_1B+a_0
\]
has the form of a polynomial evaluated at $B$. For example, if
\[
p(x)=5x^2+2x+7,
\]
then
\[
527=p(10).
\]
Thus a positional numeral can be treated as a polynomial whose coefficients are digits and whose variable is evaluated at the chosen base.

\section{The General Multiplication Formula}

Suppose two nonnegative integers are represented in base $B$ as
\[
A=\sum_{i=0}^{m}a_iB^i
\]
and
\[
C=\sum_{j=0}^{n}c_jB^j.
\]
Multiplying the sums gives
\[
AC=
\left(\sum_{i=0}^{m}a_iB^i\right)
\left(\sum_{j=0}^{n}c_jB^j\right).
\]
Using the distributive property,
\[
AC=\sum_{i=0}^{m}\sum_{j=0}^{n}a_ic_jB^{i+j}.
\]
Terms with the same total exponent can be grouped together. Define
\[
s_k=\sum_{i+j=k}a_ic_j.
\]
Then
\[
\boxed{AC=\sum_{k=0}^{m+n}s_kB^k},
\]
where
\[
\boxed{s_k=\sum_{i+j=k}a_ic_j}.
\]
Equivalently,
\[
\boxed{
\left(\sum_i a_iB^i\right)
\left(\sum_j c_jB^j\right)
=
\sum_k\left(\sum_{i+j=k}a_ic_j\right)B^k
}.
\]

The coefficient $s_k$ is obtained by multiplying every pair of digits whose positions add to $k$. This explains the diagonal pattern seen in traditional long multiplication.

\section{Multiplication as Convolution}

If the digit sequences are written from least significant to most significant,
\[
(a_0,a_1,\ldots,a_m)
\]
and
\[
(c_0,c_1,\ldots,c_n),
\]
then the sequence
\[
s_k=\sum_{i+j=k}a_ic_j
\]
is their discrete convolution.

For two two-digit numbers, the least-significant-first digit sequences are
\[
(b,a)
\qquad\text{and}\qquad
(d,c).
\]
Their convolution is
\[
(bd,\ ad+bc,\ ac),
\]
which gives
\[
bd+(ad+bc)B+acB^2.
\]
Thus long multiplication has two conceptually separate stages:
\begin{enumerate}[label=\arabic*.]
\item Convolve the digit sequences.
\item Carry to convert the resulting coefficients into valid digits.
\end{enumerate}

This same convolution structure appears in polynomial multiplication, digital signal processing, probability, and several fast multiplication algorithms.

\section{Carrying as Base-$B$ Normalization}

The convolution coefficients are not necessarily valid base-$B$ digits. For example, the raw coefficients for
\[
23\times45
\]
are
\[
(15,22,8),
\]
ordered from the $10^0$ position upward. However, decimal digits must lie between $0$ and $9$.

The expression before carrying is
\[
15(10^0)+22(10^1)+8(10^2).
\]
Normalize it from right to left:
\[
15=5+1(10).
\]
Write $5$ in the ones position and carry $1$ into the tens position. The tens coefficient becomes
\[
22+1=23.
\]
Then
\[
23=3+2(10).
\]
Write $3$ in the tens position and carry $2$ into the hundreds position. The hundreds coefficient becomes
\[
8+2=10.
\]
Finally,
\[
10=0+1(10),
\]
which produces the digits
\[
1,0,3,5.
\]
Therefore,
\[
23\times45=1035.
\]

More generally, suppose the raw coefficient in position $k$, including the incoming carry, is $t_k$. By the division algorithm, there are unique integers $q_{k+1}$ and $d_k$ such that
\[
t_k=Bq_{k+1}+d_k,
\qquad
0\le d_k<B.
\]
Equivalently,
\[
d_k=t_k\bmod B
\]
and
\[
q_{k+1}=\left\lfloor\frac{t_k}{B}\right\rfloor.
\]
The remainder $d_k$ becomes the digit in position $k$, and the quotient $q_{k+1}$ is carried into position $k+1$.

Therefore,
\[
\boxed{
\text{positional multiplication}
=
\text{convolution}
+
\text{base-}B\text{ normalization}
}.
\]

\section{Example in Binary}

Consider
\[
1011_2\times110_2.
\]
The corresponding polynomials are
\[
x^3+x+1
\]
and
\[
x^2+x.
\]
Their product is
\begin{align*}
(x^3+x+1)(x^2+x)
&=x^5+x^4+x^3+2x^2+x.
\end{align*}
Evaluating at $x=2$ gives the correct numerical value, but the coefficient $2$ is not a valid binary digit. Carrying in base two transforms the raw coefficients into
\[
1000010_2.
\]
Indeed,
\[
1011_2=11_{10},
\qquad
110_2=6_{10},
\]
and
\[
11\cdot6=66=1000010_2.
\]
The multiplication rule is unchanged when moving from decimal to binary. Only the base, the allowed digits, and the carrying threshold change.

\section{Arbitrary Bases}

For any integer base $B\ge 2$, two two-digit numerals satisfy
\[
(aB+b)(cB+d)=acB^2+(ad+bc)B+bd.
\]
Setting $B=10$ gives decimal multiplication:
\[
100ac+10(ad+bc)+bd.
\]
Setting $B=2$ gives the corresponding expression in binary:
\[
4ac+2(ad+bc)+bd.
\]
Setting $B=16$ gives
\[
256ac+16(ad+bc)+bd.
\]

The algebraic structure is independent of the base. What changes is the value assigned to each position, the permitted digit set, and the point at which carrying occurs.

\section{Connection to Faster Multiplication}

The ordinary school algorithm forms every digit-by-digit product. If each number has approximately $n$ digits, this takes on the order of $n^2$ single-digit multiplications.

Once numbers are recognized as polynomials evaluated at a base, faster polynomial-multiplication ideas can also be applied to integers. Karatsuba multiplication begins by splitting two numbers into high and low parts:
\[
X=x_1B^m+x_0,
\qquad
Y=y_1B^m+y_0.
\]
Their direct product is
\[
XY=x_1y_1B^{2m}+(x_1y_0+x_0y_1)B^m+x_0y_0.
\]
The middle term can be computed using
\[
x_1y_0+x_0y_1
=(x_1+x_0)(y_1+y_0)-x_1y_1-x_0y_0.
\]
This replaces four recursive multiplications with three. For sufficiently large integers, the savings become substantial.

Even faster approaches use evaluation, interpolation, and Fourier-transform methods to compute polynomial convolutions efficiently. The conceptual path is
\[
\text{digit multiplication}
\longrightarrow
\text{polynomial multiplication}
\longrightarrow
\text{convolution}
\longrightarrow
\text{fast integer multiplication}.
\]

\section{Educational Significance}

The polynomial interpretation creates a bridge between arithmetic and algebra. For example,
\[
(2x+3)(4x+5)=8x^2+22x+15.
\]
Evaluating at $x=10$ gives
\[
8(100)+22(10)+15=1035,
\]
which is the same as
\[
23\times45.
\]
The only additional step required in positional notation is carrying.

This viewpoint helps explain:
\begin{itemize}
\item why place values are powers of the base,
\item why partial products shift position,
\item why cross-products combine,
\item why carrying works,
\item why the same arithmetic applies in other bases,
\item why polynomial multiplication resembles long multiplication, and
\item why convolution appears in integer arithmetic.
\end{itemize}

\section{Conclusion}

Beginning with
\[
(10a+b)(10c+d)=100ac+10(ad+bc)+bd,
\]
we uncover the general structure of positional multiplication. A base-$B$ number is represented by
\[
\sum_i a_iB^i.
\]
Multiplying two such numbers gives
\[
\sum_k\left(\sum_{i+j=k}a_ic_j\right)B^k.
\]
The inner sum is a discrete convolution of the digit sequences. Because the convolution coefficients may be greater than or equal to the base, carrying is required to convert them into valid digits.

The resulting principle is
\[
\boxed{
\text{place-value arithmetic is polynomial arithmetic evaluated at the base, followed by carrying}
}.
\]

What begins as a formula for multiplying two-digit decimal numbers therefore leads naturally to arbitrary bases, polynomial algebra, convolution, and modern multiplication algorithms.

\begin{thebibliography}{9}

\bibitem{bernstein}
D. J. Bernstein,
\emph{Multidigit Multiplication for Mathematicians}, 2001.
Available at \url{https://cr.yp.to/papers/m3.pdf}.

\bibitem{howe-epp}
R. Howe and S. Epp,
\emph{Taking Place Value Seriously: Arithmetic, Estimation, and Algebra}.
Available at \url{https://static.ias.edu/pcmi/2005/documents/rhplacevalue05-1.pdf}.

\bibitem{knuth}
D. E. Knuth,
\emph{The Art of Computer Programming, Volume 2: Seminumerical Algorithms},
3rd ed., Addison-Wesley, 1997.

\end{thebibliography}

\end{document}
